% =====================================================================
\section{Introduction}\label{sec:intro}

\subsection{Motivation and client context}
The upper few metres of the ocean carry oil, plastic debris, larvae, life rafts and lost cargo. Operational search-and-rescue and spill-response systems therefore depend on short-range estimates of where surface water, and objects in it, will be in the next hours to days \citep{breivik2008,vansebille2018}. Better 24-hour position estimates matter for three kinds of user. \emph{Search and rescue:} planners derive a search area from a last known position and a drift prediction, and the area needed to contain a target grows with the square of the position uncertainty. Objects such as life rafts, however, also respond to wind (``leeway''), so drifter skill transfers to them only in part \citep{breivik2008,breivik2011}. \emph{Pollution and debris tracking:} short-range transport estimates help decide where to monitor and intervene. Oil and floating debris, however, experience windage, wave-induced Stokes drift and other processes that a drogued drifter is designed to avoid \citep{vansebille2020}. \emph{Oceanographic monitoring:} users of the GDP record need to know when a drifter's recent track is a reliable guide to its near future, and when it is not. This study evaluates none of these applications directly; it quantifies retrospective 24-hour forecast skill for GDP drifters, which is a necessary first step for all three.

Satellite-tracked surface drifters provide the most direct observational record of this motion. The NOAA Global Drifter Program (GDP) maintains an array of drogued buoys that follow near-surface currents. Their tracks underpin global estimates of surface circulation and are a reference for evaluating ocean models and calibrating satellite products \citep{lumpkin2013,centurioni2019}.

This project uses the hourly GDP product \citep{elipot2016} and its companion hourly sea-surface temperature (SST) dataset \citep{elipot2022sst}, supplied as the course dataset. The hourly product interpolates irregular Argos and GPS fixes onto a uniform one-hour grid and provides positions, velocities and uncertainty estimates. This resolution is what makes a 24-hour forecasting study with hourly inputs feasible. A question of direct interest to data providers and users is therefore how much short-range predictability can be extracted from the product's own kinematic information, and under which motion conditions simple extrapolation breaks down.

We study the Gulf Stream, a narrow, fast western boundary current whose path downstream of Cape Hatteras meanders and has become less stable in recent decades \citep{andres2016}. Strong spatial gradients and rapid changes of direction make it a demanding test region. A drifter moving at \SI{1}{m\,s^{-1}} travels about \SI{86}{km} in 24 hours, so a small error in the assumed velocity quickly becomes a large position error.

\subsection{Related work}
Data-driven forecasting of drifter motion has clear precedents. \citet{aksamit2020} combined a recurrent neural network trained on drifter motion history with a reduced Maxey--Riley model and found improvements over the physical model alone for Gulf of Mexico deployments. \citet{grossi2025} predicted Gulf of Mexico drifter trajectories over 24-hour and five-day windows. They reported that fully connected networks using velocity histories did not outperform autoregressive baselines, whereas a model sharing information between drifters performed better in many cases: added flexibility does not guarantee skill. \citet{lin2026} benchmarked kinematic and machine-learning models for ten-minute forecasts of a single buoy under explicit leakage control. Constant velocity performed best overall, and trajectory-geometry-dependent gains from learned corrections were limited. Together these studies suggest that (i) strong kinematic baselines are essential, (ii) evaluation must prevent leakage between training and test trajectories, and (iii) aggregate scores can conceal condition-specific behaviour. Our study combines these elements at a 24-hour horizon on many held-out drifters. We make no claim to introduce a new forecasting algorithm or to be the first to use motion history.

\subsection{Research question, objectives and scope}
The primary research question, fixed in the project proposal, is:
\begin{quote}\itshape
How do 24-hour position forecast errors vary with a drifter's recent speed, acceleration and turning, and under which motion conditions does trajectory history improve prediction beyond current-state information and physical baselines?
\end{quote}
The objectives were to (O1) construct comparable 24-hour forecasts from Constant Velocity, Mean-flow Advection and learned residual-correction models; (O2) quantify how median and upper-tail errors vary with pre-forecast motion characteristics; (O3) compare each learned model with and without history features to identify where history helps; and (O4) report the results reproducibly, with uncertainty that respects the dependence between forecasts from the same drifter. We extended the original plan with three further model families, spatial, temporal and SST analyses, and an SST feature ablation.

The scope is deliberately narrow: one region, one season (July--September), one forecast horizon (24\,h) and one history length (24\,h). Forecast ``difficulty'' is defined relative to the tested methods and data, not as an intrinsic limit on ocean predictability, and all regime comparisons are associations rather than causal effects.

\Cref{sec:data,sec:framework,sec:design} describe the data, forecasting framework and experimental design. \Cref{sec:overall,sec:motion,sec:context} present the results, and \cref{sec:discussion,sec:conclusion} discuss limitations, recommendations and conclusions. \Cref{sec:ai} reflects on the use of AI tools.
